\documentclass[11pt,a4paper,ngerman]{book}
\usepackage{ngerman}
\usepackage{graphicx}
\usepackage[utf8]{inputenc}
\begin{document}

% Start der Titelseite
%---------------------
\begin{titlepage}
\begin{center}
% HWR-Logo ganz nach oben
\includegraphics[width=0.45\textwidth]{./hwr-logo}\\[2cm]

\textsc{\Large Schriften zur Wirtschaftsinformatik}
\\[2cm]

{\Huge Konzepte der Web-Programmierung am Beispiel konkreter Softwareentwicklungsprojekte} 
\\[2cm]

% Autoren und Editor links- bzw. rechtsbündig einsetzen
\begin{minipage}{0.4\textwidth}
\begin{flushleft} \large
\emph{Autoren:} \\
Studierende des \\ 
des dualen Studiengangs Wirtschaftsinformatik 
\end{flushleft}
\end{minipage}
\begin{minipage}{0.4\textwidth}
\begin{flushright} \large
\emph{Autor und Editor:}\\
Prof. Dr. \\ Gert Faustmann
\end{flushright}
\end{minipage}

% Auffüllen und Jahr drucken
\vfill
2015
\end{center}
\end{titlepage}
%--------------------
% Ende der Titelseite


\chapter{Einleitung}

\chapter{Grundlagen der Web-Programmierung}
\section{Ansätze}
Browser, Server, HTML, Skripte, Server/Clientseitige Verarbeitung, ...

\section{Vorteile}
Scaffolding, DRY, ...

\section{Herausforderungen}
Vielfalt an Technologien, Wartbarkeit, Effizienz,...



\chapter{Beschreibung der Projekte}

% Rahmenbedingungen noch von GF
\section{Rahmenbedingungen}
Zeitrahmen, formaler Rahmen, inhaltliche Vorgaben, Vorwissen (Ziele WI-Studium, Aufbau des Studiums, bisherige Programmierveranstaltungen


% Bitte nur ab diesem Abschnitt Ergebnisse aus den Projekten einfügen.
% Überschrift "Projekt 1" durch eigenen Titel ersetzen.
\section{Projekt 1}

\subsection{Kurzdarstellung}
Ziele des Systems, Zielgruppe

\subsection{Projektrahmen}
Mitglieder(-zahl), Vorkenntnisse (anonym?)

\subsection{Domänenanalyse}
\subsubsection{Funktionale Anforderungen}
\subsubsection{Informationsmodell}
\subsubsection{Nutzerschnittstelle}

\subsection{Technische Plattform}
\subsubsection{Architektur}
\subsubsection{Verwendete Komponenten}

\subsection{Entwicklungsprozess}
Ablauf, Arbeitsteilung, zeitliche Ressourcen

\subsection{Bewertung}
\subsubsection{Vorteile}
Entwicklungszeit, Integrationsfähigkeit,...
\subsubsection{Nachteile}
Komplexität, Einarbeitungszeit, Schwierigkeit der Integration, ...

\end{document}
